\section{Native moments and calibrated posterior risk}
\begin{lemma}[Uniform native moments]
\label{lem:moments}
There is $C=C_{H,A,\beta}$ such that, for every input and every scale~\eqref{eq:scales},
\[
 \E_\mu |z_w(x)|^4\le C,\qquad
 \E_\mu R(z_w)^2\le C,\qquad
 \E_\mu\whR_S(z_w)\le C.
\]
Consequently $Z_\lambda\ge e^{-C/(\lambda\sigma^2)}$ for every sample.
\end{lemma}
\begin{proof}
The first-layer preactivation variance is at most $2$. The inequality
$|\phi(t)|\le A|t|+\beta$, Gaussian fourth moments and
$(N^{-1}\sum_i h_i^2)^2\le N^{-1}\sum_i h_i^4$ inductively bound
$\E(N^{-1}\|h^\ell\|^2)^2$ by a depth-dependent constant. Conditional on the final hidden layer, the readout is Gaussian with variance
$\sigma^2\|h^H\|^2/N$. Since $\sigma^2\le1$, its fourth moment has the claimed bound. Jensen and the bounded labels give the risk moment statements. Jensen applied to the exponential defining $Z_\lambda$ gives its lower bound.
\end{proof}

